\documentclass[10pt,a4paper]{amsart}
\usepackage[margin=19mm]{geometry}
\usepackage{amsmath,amssymb,mathtools}
\usepackage[scr=rsfs,cal=euler]{mathalfa}
\usepackage{xfrac}
\usepackage[unicode,hidelinks,bookmarks=false]{hyperref}
\newcommand{\ie}{i.e.\ }
\newcommand{\cf}{cf.\ }
\newcommand{\resp}{respectively\ }
\newcommand{\Subsection}{\subsection}
\providecommand{\nobreakdash}{\nobreak\hbox{-}}
\setlength{\emergencystretch}{2em}
\multlinegap0pt
\usepackage{amssymb}
\usepackage{mathtools}
\usepackage[shortlabels]{enumitem}
\newtheorem{proposition}{Proposition}
\title{Revisiting the Behrens--Fisher Problem: Validity-First Optimality}
\author{Xiao Wang, Chuanhai Liu}
\date{Source: arXiv:2606.07847v1 (CC BY 4.0)}
\begin{document}
\maketitle

\noindent\emph{Source and licence.} Revisiting the Behrens--Fisher Problem: Validity-First Optimality. Version arXiv:2606.07847v1 (CC BY 4.0), distributed by arXiv under the Creative Commons Attribution 4.0 International licence, http://creativecommons.org/licenses/by/4.0/, as recorded in the arXiv licence metadata for this exact version and shown on the abstract page https://arxiv.org/abs/2606.07847v1. Full licence notice: \texttt{licenses/arxiv-2606.07847-CC-BY-4.0.txt}.\par\smallskip

\noindent\textit{Editorial scope.} Counted unit: the article's proposition prop:nonregular (source0.tex lines 428--434). The article's complete original proof of it is delivered with it (source0.tex lines 436--488, one \texttt{proof} environment of 53 lines). No mathematics is written, repaired or re-derived: the statement and the proof are the article's own lines, in the article's own notation, and no retained line is substituted or re-flowed.  Retained uncounted as the source-local context the proof needs: lines 339--342; lines 397--408.  Nothing was omitted from any retained span, so the package carries no omission record.  No retained line contains a citation, so the wrapper carries no bibliography and no bibliography entry.  No retained line contains a figure environment, an image-inclusion command or drawing code, so no image is delivered and the package declares no asset dependency.  Editorial formatting only, in addition: an AMS \texttt{amsart} preamble that restates the article's own declarations and macros used by the retained lines, namely the environments \texttt{proposition} on separate counters, together with the standard packages those lines need.  The retained environments print in this document as Proposition 1 in order of appearance, whereas the article prints the same items with the numbering of its own full document.  The complete native source of the article as distributed in the arXiv e-print for this exact version is bundled unchanged as \texttt{sources/202280/source0.tex}.
\begin{align}
\bar p(X,\psi)&=\bar a(V_1,\psi), \label{eq:regular-interest}\\
c(X,V_2,\psi,\xi)&=0, \label{eq:regular-nuisance}
\end{align}

\begin{align}
T_\psi
&=
\frac{\bar Y-\psi}{f(S_1,S_2)}
=
Z_1(\xi), \label{eq:bf-first}\\
R
&=
\frac{n_1S_2^2}{n_2S_1^2}
=
Z_2\frac{1-\xi}{\xi}, \label{eq:bf-second}
\end{align}

\begin{proposition}[Nonregularity]
\label{prop:nonregular}
The reduced Behrens--Fisher association \eqref{eq:bf-first}--\eqref{eq:bf-second}
is not regular in the IM sense of Martin and Liu: it cannot be brought to the
form \eqref{eq:regular-interest}--\eqref{eq:regular-nuisance} with a
nuisance-free interest equation.
\end{proposition}

\begin{proof}
Regularity in the sense of
\eqref{eq:regular-interest}--\eqref{eq:regular-nuisance} would require an
auxiliary coordinate for \(\psi\) whose distribution does not depend on the
nuisance, together with a companion equation that contributes no information
about \(\psi\).  We show that neither requirement can be met.

\emph{Step 1: the interest auxiliary is nuisance-dependent.}
Fix \(\xi\in(0,1)\) and write \(Q_k=(n_k-1)U_{2k}^2\sim\chi^2_{\nu_k}\) with
\(\nu_k=n_k-1\), independent of \(U_1=N\sim N(0,1)\).  Then \(U_{2k}^2=Q_k/\nu_k\)
and
\[
Z_1(\xi)
=
\frac{U_1}{\{\xi U_{21}^2+(1-\xi)U_{22}^2\}^{1/2}}
\overset{d}{=}
\frac{N}{\{\xi Q_1/\nu_1+(1-\xi)Q_2/\nu_2\}^{1/2}} .
\]
Thus, conditionally on \((Q_1,Q_2)\), \(Z_1(\xi)\) is centered normal with a
random scale, so its marginal law is a scale mixture of \(N(0,1)\) whose mixing
law is that of the denominator.  This mixing law varies with \(\xi\), and the
variation is genuine rather than removable by a change of auxiliary variable.
To exhibit this, it suffices to display two values of \(\xi\) giving distinct
laws.  As \(\xi\to1\) the denominator converges to \((Q_1/\nu_1)^{1/2}\), so
\(Z_1(\xi)\Rightarrow t_{\nu_1}\); as \(\xi\to0\) it converges to
\((Q_2/\nu_2)^{1/2}\), so \(Z_1(\xi)\Rightarrow t_{\nu_2}\).  When
\(\nu_1\neq\nu_2\) these limits already differ.  When \(\nu_1=\nu_2=\nu\), take
instead \(\xi=\tfrac12\): there \(\xi U_{21}^2+(1-\xi)U_{22}^2=(Q_1+Q_2)/(2\nu)\)
with \(Q_1+Q_2\sim\chi^2_{2\nu}\), whence \(Z_1(\tfrac12)\sim t_{2\nu}\neq
t_{\nu}\).  In every case the family \(\{\,\text{law of }Z_1(\xi):\xi\in(0,1)\,\}\)
is non-degenerate, so the interest auxiliary is not nuisance-free and
\eqref{eq:bf-first} cannot play the role of the regular interest equation
\eqref{eq:regular-interest}.

\emph{Step 2: the companion equation does not restore regularity.}
Regularity could still hold if \eqref{eq:bf-second} determined \(\xi\) from
observables alone, thereby fixing the law in Step~1.  Solving
\eqref{eq:bf-second} for \(\xi\) gives
\[
\xi=\frac{Z_2}{Z_2+R},
\qquad
R=\frac{n_1S_2^2}{n_2S_1^2}\ \text{observed},\quad
Z_2=\frac{U_{22}^2}{U_{21}^2}\ \text{unobserved}.
\]
Hence \(\xi\) is not pinned down by the data; it is recovered only through the
unobserved auxiliary \(Z_2\), which must itself be predicted.  Eliminating the
nuisance parameter merely trades it for a second auxiliary coordinate.

The two steps together show that no rewriting of the form
\eqref{eq:regular-interest}--\eqref{eq:regular-nuisance} with a nuisance-free
interest equation exists.  The association therefore remains genuinely
two-dimensional after all regular marginalization has been exhausted.
\end{proof}

\end{document}
